% ============================================================
%  PEMBAHASAN SUPER DETAIL — LATIHAN SOAL & STUDI KASUS
%  BAB 1: EKSPONEN DAN LOGARITMA
%  Kurikulum Merdeka — Fase E (Kelas 10)
%  Kompilasi: pdfLaTeX / XeLaTeX (disarankan Overleaf)
%  Catatan: Soal Akhir Bab TIDAK dibahas di sini (sesuai permintaan).
% ============================================================
\documentclass[11pt,a4paper]{report}

% ---------- PAKET ----------
\usepackage[utf8]{inputenc}
\usepackage[T1]{fontenc}
\usepackage[bahasa]{babel}
\usepackage{amsmath,amssymb,amsthm,mathtools}
\usepackage{geometry}
\usepackage{xcolor}
\usepackage{graphicx}
\usepackage{enumitem}
\usepackage{tikz}
\usetikzlibrary{calc,arrows.meta,angles,quotes,positioning,decorations.pathreplacing}
\usepackage{pgfplots}
\pgfplotsset{compat=1.18}
\usepackage[most]{tcolorbox}
\usepackage{fancyhdr}
\usepackage{titlesec}
\usepackage{array,booktabs}
\usepackage{multicol}
\usepackage{pifont}
\usepackage[colorlinks=true,linkcolor=biru,urlcolor=biru]{hyperref}

\geometry{a4paper,margin=2.5cm,headheight=15pt}

% ---------- WARNA ----------
\definecolor{biru}{HTML}{1F4E79}
\definecolor{birumuda}{HTML}{D6E4F0}
\definecolor{hijau}{HTML}{2E7D32}
\definecolor{hijaumuda}{HTML}{E3F2E1}
\definecolor{oranye}{HTML}{C75B12}
\definecolor{oranyemuda}{HTML}{FBE8DA}
\definecolor{abu}{HTML}{F2F2F2}
\definecolor{ungu}{HTML}{6A1B9A}
\definecolor{ungumuda}{HTML}{EDE2F3}
\definecolor{merah}{HTML}{C62828}
\definecolor{merahmuda}{HTML}{FCE4E4}
\definecolor{tosca}{HTML}{00838F}
\definecolor{toscamuda}{HTML}{E0F2F1}
\definecolor{emas}{HTML}{8D6E00}
\definecolor{emasmuda}{HTML}{FFF6D6}

% ---------- HEADER / FOOTER ----------
\pagestyle{fancy}
\fancyhf{}
\renewcommand{\headrulewidth}{0.8pt}
\renewcommand{\footrulewidth}{0pt}
\fancyhead[L]{\small\textcolor{ungu}{\leftmark}}
\fancyhead[R]{\small\textcolor{ungu}{Pembahasan Latihan — Bab 1}}
\fancyfoot[C]{\thepage}

% ---------- GAYA JUDUL BAB ----------
\titleformat{\chapter}[display]
  {\normalfont\huge\bfseries\color{ungu}}
  {\filright\large PEMBAHASAN BAB \thechapter}{8pt}{\Huge}
\titlespacing*{\chapter}{0pt}{0pt}{20pt}

% ---------- LINGKUNGAN TEOREMA ----------
\newtheoremstyle{gaya}{6pt}{6pt}{}{}{\bfseries\color{biru}}{.}{.5em}{}
\theoremstyle{gaya}
\newtheorem{definisi}{Definisi}[chapter]
\newtheorem{sifat}{Sifat}[chapter]

% ---------- KOTAK: SOAL (ungu) ----------
\newtcolorbox{soalbox}[1]{
  enhanced, breakable, colback=ungumuda, colframe=ungu, boxrule=0.8pt,
  arc=3pt, left=10pt, right=10pt, top=8pt, bottom=8pt,
  title={\textbf{Soal #1}}, fonttitle=\bfseries\color{white}, coltitle=white,
  attach boxed title to top left={xshift=10pt,yshift=-10pt},
  boxed title style={colback=ungu,arc=2pt}}

% ---------- KOTAK: KONSEP KUNCI (hijau) ----------
\newtcolorbox{ingat}[1]{
  enhanced, breakable, colback=hijaumuda, colframe=hijau, boxrule=0.6pt,
  arc=3pt, left=10pt, right=10pt, top=6pt, bottom=6pt,
  title={\textbf{Konsep Kunci: #1}}, fonttitle=\bfseries\color{white}, coltitle=white,
  attach boxed title to top left={xshift=10pt,yshift=-10pt},
  boxed title style={colback=hijau,arc=2pt}}

% ---------- KOTAK: JEBAKAN (merah) ----------
\newtcolorbox{jebakan}{
  enhanced, breakable, colback=merahmuda, colframe=merah, boxrule=0.6pt,
  arc=3pt, left=10pt, right=10pt, top=6pt, bottom=6pt,
  title={\textbf{\ding{55}\ Jebakan \& Kesalahan yang Sering Terjadi}},
  fonttitle=\bfseries\color{white}, coltitle=white,
  attach boxed title to top left={xshift=10pt,yshift=-10pt},
  boxed title style={colback=merah,arc=2pt}}

% ---------- KOTAK: VARIASI SOAL (oranye) ----------
\newtcolorbox{variasi}{
  enhanced, breakable, colback=oranyemuda, colframe=oranye, boxrule=0.6pt,
  arc=3pt, left=10pt, right=10pt, top=6pt, bottom=6pt,
  title={\textbf{\ding{72}\ Bentuk Modifikasi Soal Ini}},
  fonttitle=\bfseries\color{white}, coltitle=white,
  attach boxed title to top left={xshift=10pt,yshift=-10pt},
  boxed title style={colback=oranye,arc=2pt}}

% ---------- KOTAK: STUDI KASUS (biru) ----------
\newtcolorbox{studikasus}[1]{
  enhanced, breakable, colback=abu, colframe=biru, boxrule=0.8pt,
  arc=3pt, left=10pt, right=10pt, top=8pt, bottom=8pt,
  title={\textbf{Studi Kasus: #1}}, fonttitle=\bfseries\color{white}, coltitle=white,
  attach boxed title to top left={xshift=10pt,yshift=-10pt},
  boxed title style={colback=biru,arc=2pt}}

% ---------- KOTAK: TIPS (tosca) ----------
\newtcolorbox{tips}{
  enhanced, breakable, colback=toscamuda, colframe=tosca, boxrule=0.6pt,
  arc=3pt, left=10pt, right=10pt, top=6pt, bottom=6pt,
  borderline west={3pt}{0pt}{tosca}}

% ---------- PERINTAH BANTU ----------
\newcommand{\jawab}{\par\smallskip\noindent\textit{\textcolor{oranye}{Penyelesaian langkah demi langkah:}}\par\smallskip}
\newcommand{\langkah}[1]{\par\smallskip\noindent\textbf{\textcolor{biru}{Langkah #1.}}\ }
\newcommand{\jadi}{\par\smallskip\noindent$\Rightarrow$\ \textbf{Jadi, }}
% Judul tiap nomor soal
\newcommand{\soalno}[2]{%
  \par\bigskip
  \noindent{\large\bfseries\color{ungu}\rule[-2pt]{4pt}{14pt}\ Soal #1}%
  \ {\normalfont\itshape\color{gray}(#2)}\par\nopagebreak\smallskip}

% ============================================================
\begin{document}
\setcounter{chapter}{1}

% ---------- HALAMAN JUDUL ----------
\begin{titlepage}
\centering
\vspace*{2cm}
{\color{ungu}\rule{\textwidth}{2pt}}\\[0.6cm]
{\Huge\bfseries\color{ungu} PEMBAHASAN\\[0.3cm] LATIHAN SOAL}\\[0.4cm]
{\Large BAB 1 — Eksponen dan Logaritma}\\[0.2cm]
{\color{ungu}\rule{\textwidth}{2pt}}\\[1.2cm]

{\large Matematika SMA — Fase E (Kelas 10)}\\[0.2cm]
{\large Kurikulum Merdeka}\\[2cm]

\begin{tcolorbox}[enhanced,width=0.85\textwidth,colback=ungumuda,colframe=ungu,
  arc=4pt,boxrule=1pt]
\centering\small
Buku pembahasan ini mengupas \textbf{setiap soal pada kotak ungu (Latihan Soal)}
di tiap subbab, \emph{selangkah demi selangkah seperti mengajari dari nol}:
konsep yang dipakai, jebakan yang sering menjerat, cara soal bisa dimodifikasi,
dan visualisasi. \emph{Soal Akhir Bab tidak dibahas di sini.}
\end{tcolorbox}

\vfill
{\large Tahun Pelajaran 2026/2027}
\end{titlepage}

\tableofcontents
\thispagestyle{empty}
\clearpage

% ============================================================
\chapter*{Pembahasan Bab 1: Eksponen dan Logaritma}
\addcontentsline{toc}{chapter}{Pembahasan Bab 1: Eksponen dan Logaritma}
\markboth{PEMBAHASAN BAB 1}{}
% ============================================================

\begin{tips}
\textbf{Cara membaca buku ini.} Setiap soal dibahas dengan pola tetap:
\textbf{(1) soal ditulis ulang}, \textbf{(2) konsep kunci} yang dipakai,
\textbf{(3) penyelesaian langkah demi langkah}, \textbf{(4) jebakan} yang
sering bikin salah, dan \textbf{(5) variasi} bentuk soal yang mungkin muncul di
ujian. Jangan lewati kotak merah (jebakan) --- di situlah nilai sering hilang.
\end{tips}

% ============================================================
\section*{Studi Kasus: Pertumbuhan Subscriber Kanal YouTube}
\addcontentsline{toc}{section}{Studi Kasus: Pertumbuhan Subscriber YouTube}
% ============================================================

\begin{studikasus}{Pertumbuhan Subscriber Kanal YouTube}
Seorang kreator memulai kanal dengan $1000$ subscriber. Setiap bulan jumlah
subscriber-nya menjadi $1{,}5$ kali lipat bulan sebelumnya.

\smallskip
\begin{center}
\begin{tabular}{c|c|c}
\toprule
Bulan ke-$n$ & Cara hitung & Subscriber \\
\midrule
0 & $1000$ & $1000$\\
1 & $1000 \times 1{,}5$ & $1500$\\
2 & $1000 \times 1{,}5 \times 1{,}5$ & $2250$\\
3 & $1000 \times 1{,}5^3$ & $3375$\\
$n$ & $1000 \times 1{,}5^{\,n}$ & ? \\
\bottomrule
\end{tabular}
\end{center}

\smallskip
\textbf{Pertanyaan.} (a) Apa pola yang muncul di kolom ``cara hitung''?
(b) Bagaimana menuliskan jumlah subscriber pada bulan ke-$n$ secara ringkas?
(c) Pada bulan ke berapa subscriber menembus $100.000$?
\end{studikasus}

\subsection*{Jawaban lengkap studi kasus}

\begin{ingat}{Perkalian berulang $\to$ pangkat}
Sama seperti $\underbrace{3+3+3+3}_{4\text{ kali}} = 4\times 3$ (penjumlahan berulang
ditulis singkat dengan perkalian), maka
$\underbrace{1{,}5\times 1{,}5\times \cdots \times 1{,}5}_{n\text{ faktor}} = 1{,}5^{\,n}$
(perkalian berulang ditulis singkat dengan \emph{pangkat}). Pangkat adalah
``mesin pelipatganda''.
\end{ingat}

\textbf{(a) Pola kolom ``cara hitung''.} Perhatikan: setiap turun satu baris,
muncul \emph{satu faktor $1{,}5$ tambahan}. Pada bulan ke-$n$, angka $1{,}5$
dikalikan sebanyak $n$ kali. Bilangan awal $1000$ tidak pernah berubah --- ia
hanya ``penumpang tetap''. Maka polanya:
\[
\text{subscriber bulan ke-}n = 1000 \times \underbrace{1{,}5\times1{,}5\times\cdots\times1{,}5}_{n\text{ faktor}}.
\]

\textbf{(b) Rumus ringkas.} Perkalian berulang itu kita padatkan menjadi pangkat:
\[
\boxed{\,U_n = 1000 \times 1{,}5^{\,n}\,}
\]
dengan $U_n$ = jumlah subscriber pada bulan ke-$n$, $1000$ = nilai awal, dan
$1{,}5$ = \emph{faktor pengali} tiap bulan.

\textbf{(c) Kapan menembus $100.000$?} Kita ingin mencari $n$ dari
\[
1000\times 1{,}5^{\,n} = 100.000 \;\Longrightarrow\; 1{,}5^{\,n} = 100.
\]
Di sinilah letak masalahnya: \textbf{$n$ terjebak di posisi pangkat}. Perkalian
biasa tidak bisa ``menariknya turun''. Alat yang bisa melakukannya adalah
\emph{logaritma} (dibahas di Subbab 1.3):
\[
n = {}^{1{,}5}\!\log 100 = \frac{\log 100}{\log 1{,}5} = \frac{2}{0{,}176\ldots} \approx 11{,}4.
\]
Karena bulan haruslah bilangan bulat dan pada $n=11$ belum tembus, maka
\jadi subscriber pertama kali melampaui $100.000$ pada \textbf{bulan ke-12}.

\begin{tips}
\textbf{Pesan besar bab ini:} \emph{eksponen} memproyeksikan pertumbuhan ke depan
(``berapa banyak nanti?''), sedangkan \emph{logaritma} membongkar waktunya
(``kapan tercapai?''). Keduanya saling membalik.
\end{tips}

% ############################################################
\section*{Subbab 1.1 — Bilangan Berpangkat (Eksponen)}
\addcontentsline{toc}{section}{Subbab 1.1 — Bilangan Berpangkat}
% ############################################################

\begin{ingat}{Delapan sifat eksponen (dipakai di seluruh subbab ini)}
Untuk $a,b\neq 0$ dan $m,n$ rasional:
\[
\begin{array}{lll}
(1)\ a^m\cdot a^n=a^{m+n} & (2)\ \dfrac{a^m}{a^n}=a^{m-n} & (3)\ (a^m)^n=a^{mn}\\[8pt]
(4)\ (ab)^n=a^n b^n & (5)\ \left(\dfrac ab\right)^n=\dfrac{a^n}{b^n} & (6)\ a^0=1\\[8pt]
(7)\ a^{-n}=\dfrac{1}{a^n} & (8)\ a^{m/n}=\sqrt[n]{a^m} &
\end{array}
\]
\textbf{Inti:} kali $\to$ pangkat \emph{ditambah}; bagi $\to$ pangkat
\emph{dikurang}; pangkat-dari-pangkat $\to$ \emph{dikali}.
\end{ingat}

% ---------------- SOAL 1.1.1 ----------------
\soalno{1.1.1}{sifat kali \& bagi pangkat}
\begin{soalbox}{}
Sederhanakan $\dfrac{a^4 b^{-3}}{a^{-2} b}$.
\end{soalbox}

\jawab
\langkah{1} Kelompokkan menurut basis yang sama. Basis $a$ diurus sendiri, basis
$b$ diurus sendiri --- jangan dicampur.
\[
\frac{a^4 b^{-3}}{a^{-2} b^{1}} = \underbrace{\frac{a^4}{a^{-2}}}_{\text{basis }a} \cdot \underbrace{\frac{b^{-3}}{b^{1}}}_{\text{basis }b}.
\]
\langkah{2} Terapkan sifat (2) $\dfrac{a^m}{a^n}=a^{m-n}$. Kuncinya: pangkat atas
\emph{dikurangi} pangkat bawah. Hati-hati tanda minus:
\[
a^{4-(-2)} = a^{4+2} = a^{6}, \qquad b^{-3-1} = b^{-4}.
\]
\langkah{3} Ubah pangkat negatif menjadi pecahan dengan sifat (7)
$a^{-n}=\frac{1}{a^n}$:
\[
a^6 b^{-4} = \frac{a^6}{b^4}.
\]
\jadi $\dfrac{a^4 b^{-3}}{a^{-2}b} = \dfrac{a^6}{b^4}$.

\begin{jebakan}
\begin{itemize}[leftmargin=*,topsep=2pt]
  \item \textbf{Salah tanda pada $4-(-2)$.} Banyak yang menulis $a^{4-2}=a^2$.
  Ingat: mengurangi bilangan negatif = \emph{menambah}. Bayangkan $-(-2)$ seperti
  ``musuh dari musuh adalah teman''.
  \item \textbf{Mencampur basis.} Menulis $\frac{a^4 b^{-3}}{a^{-2}b}$ langsung
  jadi ``$(ab)^{\text{sesuatu}}$'' itu keliru --- $a$ dan $b$ adalah dua basis
  \emph{berbeda}, kurangi pangkatnya masing-masing.
  \item \textbf{Berhenti di $b^{-4}$.} Kalau soal minta ``sederhanakan'',
  biasanya jawaban akhir tidak boleh menyisakan pangkat negatif.
\end{itemize}
\end{jebakan}

\begin{variasi}
Soal ini bisa muncul dengan wajah berbeda tetapi ide sama:
\begin{itemize}[leftmargin=*,topsep=2pt]
  \item Diberi nilai: ``Jika $a=2,b=3$, hitung $\frac{a^4b^{-3}}{a^{-2}b}$.''
  Sederhanakan dulu jadi $\frac{a^6}{b^4}$, baru masukkan angka $=\frac{64}{81}$.
  \item Ditambah basis ketiga, mis.\ $\dfrac{a^4 b^{-3} c^2}{a^{-2} b\, c^{-1}}$
  --- perlakukan $c$ persis sama.
  \item Dibalik jadi ``buktikan $\frac{a^4b^{-3}}{a^{-2}b}=\frac{a^6}{b^4}$'',
  yang menuntut langkah yang sama.
\end{itemize}
\end{variasi}

% ---------------- SOAL 1.1.2 ----------------
\soalno{1.1.2}{pangkat pecahan = akar}
\begin{soalbox}{}
Hitung $16^{3/4}$ dan $8^{-2/3}$.
\end{soalbox}

\begin{ingat}{Membaca pangkat pecahan}
$a^{m/n} = \big(\sqrt[n]{a}\big)^m = \sqrt[n]{a^m}$. Bacalah dari
\emph{penyebut ke pembilang}: penyebut $n$ = ``akar ke-$n$'', pembilang $m$ =
``dipangkatkan $m$''. Lebih mudah \textbf{akar dulu, baru pangkat}, supaya
angkanya kecil.
\end{ingat}

\jawab
\textbf{Bagian pertama: $16^{3/4}$.}
\langkah{1} Tulis basis sebagai pangkat bilangan prima: $16 = 2^4$.
\langkah{2} Gunakan sifat (3) $(a^m)^n=a^{mn}$:
\[
16^{3/4} = (2^4)^{3/4} = 2^{4\cdot \frac34} = 2^{3} = 8.
\]
\emph{Alternatif akar dulu:} $16^{3/4} = \big(\sqrt[4]{16}\big)^3 = 2^3 = 8$. Sama.

\textbf{Bagian kedua: $8^{-2/3}$.}
\langkah{1} Pangkat negatif $\Rightarrow$ balik jadi pecahan lebih dulu (sifat 7):
$8^{-2/3} = \dfrac{1}{8^{2/3}}$.
\langkah{2} Hitung $8^{2/3}$ dengan $8=2^3$: $8^{2/3}=(2^3)^{2/3}=2^2=4$.
\langkah{3} Maka $8^{-2/3} = \dfrac{1}{4}$.
\jadi $16^{3/4}=8$ dan $8^{-2/3}=\dfrac14$.

\begin{jebakan}
\begin{itemize}[leftmargin=*,topsep=2pt]
  \item \textbf{Menganggap pangkat negatif membuat hasil negatif.}
  $8^{-2/3}$ \emph{bukan} $-4$. Tanda minus di pangkat berarti ``kebalikan''
  ($\frac1{\cdots}$), bukan lawan bilangan.
  \item \textbf{Salah arah akar/pangkat.} $16^{3/4}\neq \sqrt[3]{16^4}$.
  Penyebut $4$ = akarnya, pembilang $3$ = pangkatnya.
  \item \textbf{Menghitung $16^3$ dulu} ($=4096$) lalu akar-4 --- boleh, tapi
  angkanya besar dan rawan salah. Selalu \emph{akar dulu}.
\end{itemize}
\end{jebakan}

\begin{variasi}
$25^{3/2}$, $27^{-2/3}$, $32^{4/5}$, atau $\left(\tfrac{8}{27}\right)^{-2/3}$
(pakai sifat 5: pangkatkan pembilang dan penyebut, lalu balik). Semua memakai
trik yang sama: \emph{ubah basis ke bentuk prima berpangkat}.
\end{variasi}

% ---------------- SOAL 1.1.3 ----------------
\soalno{1.1.3}{notasi ilmiah}
\begin{soalbox}{}
Tulis $0{,}000\,000\,7$ dalam notasi ilmiah.
\end{soalbox}

\begin{ingat}{Notasi ilmiah (baku)}
Bentuknya $a\times 10^{n}$ dengan $1\le a < 10$ ($a$ satu angka sebelum koma) dan
$n$ bilangan bulat. Geser koma: ke \emph{kanan} $\Rightarrow$ pangkat
\emph{negatif} (bilangan kecil); ke \emph{kiri} $\Rightarrow$ pangkat
\emph{positif} (bilangan besar).
\end{ingat}

\jawab
\langkah{1} Cari posisi koma agar tepat satu angka tak-nol di depan koma. Angka
tak-nol pertama adalah $7$. Kita ingin ``$7{,}0$''.
\langkah{2} Hitung berapa kali koma digeser. Dari $0{,}0000007$ ke $7{,}0$ koma
bergeser $7$ langkah ke \emph{kanan}.
\[
0{,}\underset{1}{0}\underset{2}{0}\underset{3}{0}\underset{4}{0}\underset{5}{0}\underset{6}{0}\underset{7}{7}
\;\longrightarrow\; 7{,}0
\]
\langkah{3} Geser ke kanan $\Rightarrow$ pangkat negatif $\Rightarrow$ eksponen
$-7$:
\jadi $0{,}0000007 = 7\times 10^{-7}$.

\begin{jebakan}
\begin{itemize}[leftmargin=*,topsep=2pt]
  \item \textbf{Salah tanda pangkat.} Bilangan \emph{kecil} ($<1$) selalu pangkat
  \emph{negatif}. Kalau kamu dapat $7\times10^{7}$ (positif), itu berarti
  $70.000.000$ --- jelas keliru.
  \item \textbf{Salah hitung nol.} Hitung langkah geseran koma, bukan banyak nol.
  \item \textbf{$a$ tidak baku.} $70\times10^{-8}$ dan $0{,}7\times10^{-6}$ punya
  nilai sama tapi \emph{bukan} notasi ilmiah baku, karena $a$ harus $1\le a<10$.
\end{itemize}
\end{jebakan}

\begin{variasi}
Arah sebaliknya: ``Tulis $6\,500\,000$ dalam notasi ilmiah'' $\to 6{,}5\times10^6$.
Atau operasi: ``$(3\times10^{-7})(2\times10^{4})$'' $\to$ kalikan angka depan
($3\cdot2=6$) dan jumlahkan pangkat ($-7+4=-3$) $\to 6\times10^{-3}$.
\end{variasi}

% ---------------- SOAL 1.1.4 ----------------
\soalno{1.1.4}{gabungan sifat (3),(4),(2)}
\begin{soalbox}{}
Sederhanakan $\dfrac{(3x^2 y)^3}{9x^4 y^{-1}}$.
\end{soalbox}

\jawab
\langkah{1} Bongkar pembilang dengan sifat (4) $(ab)^n=a^n b^n$ dan sifat (3).
\emph{Ingat angka $3$ juga ikut dipangkatkan!}
\[
(3x^2 y)^3 = 3^3 \cdot (x^2)^3 \cdot y^3 = 27 x^{6} y^{3}.
\]
\langkah{2} Tulis ulang pecahan:
\[
\frac{27 x^6 y^3}{9 x^4 y^{-1}}.
\]
\langkah{3} Bagi bagian per bagian: angka $\frac{27}{9}=3$; basis $x$:
$x^{6-4}=x^2$; basis $y$: $y^{3-(-1)}=y^{4}$.
\jadi $\dfrac{(3x^2y)^3}{9x^4y^{-1}} = 3x^2 y^4$.

\begin{jebakan}
\begin{itemize}[leftmargin=*,topsep=2pt]
  \item \textbf{Lupa memangkatkan koefisien.} Kesalahan paling umum:
  $(3x^2y)^3 = 3x^6y^3$ (angka $3$ tidak dipangkatkan). Yang benar $3^3=27$.
  \item \textbf{Salah tanda pada $y^{3-(-1)}$}, ditulis $y^{3-1}=y^2$. Padahal
  $3-(-1)=4$.
  \item \textbf{Membagi angka dengan mengurangi} ($27-9=18$). Angka dibagi biasa
  ($27\div9=3$); yang \emph{dikurangi} hanya pangkat pada basis yang sama.
\end{itemize}
\end{jebakan}

\begin{variasi}
$\dfrac{(2a^3b^{-2})^2}{4a^{-1}b^3}$ (contoh di modul), atau dibuat
bertingkat: $\left(\dfrac{(3x^2y)^3}{9x^4y^{-1}}\right)^{-2}$. Sederhanakan
dalam kurung dulu ($=3x^2y^4$), lalu pangkatkan $-2$: $\frac{1}{9x^4y^8}$.
\end{variasi}

% ---------------- SOAL 1.1.5 ----------------
\soalno{1.1.5}{memecah pangkat penjumlahan}
\begin{soalbox}{}
Jika $2^x = 5$, tentukan nilai $2^{x+3}$.
\end{soalbox}

\begin{ingat}{Pisahkan pangkat yang dijumlah}
Sifat (1) dibaca terbalik: $a^{m+n} = a^m\cdot a^n$. Jadi ``pangkat yang berupa
penjumlahan'' bisa dipecah menjadi \emph{perkalian dua pangkat}.
\end{ingat}

\jawab
\langkah{1} Pecah $2^{x+3}$ menggunakan sifat (1):
\[
2^{x+3} = 2^{x}\cdot 2^{3}.
\]
\langkah{2} Substitusi yang diketahui: $2^x = 5$ dan $2^3=8$:
\[
2^{x+3} = 5\times 8 = 40.
\]
\jadi $2^{x+3} = 40$. (Kita \emph{tidak perlu} tahu nilai $x$ itu sendiri!)

\begin{jebakan}
\begin{itemize}[leftmargin=*,topsep=2pt]
  \item \textbf{Menambah, bukan mengali.} Menulis $2^{x+3}=2^x + 2^3 = 5+8=13$.
  \emph{Salah total.} Pangkat yang dijumlah $\to$ basis \emph{dikali}, bukan
  ditambah.
  \item \textbf{Memaksa cari $x$.} Sebagian siswa menghitung $x=\log_2 5\approx2{,}32$
  lalu $2^{5{,}32}$ --- panjang dan rawan. Cukup pecah pangkatnya.
\end{itemize}
\end{jebakan}

\begin{variasi}
``Jika $2^x=5$, tentukan $2^{x-1}$'' ($=\frac52$), atau $2^{2x}=(2^x)^2=25$, atau
$8^x=(2^3)^x=(2^x)^3=125$. Semua bertumpu pada memecah/menyusun ulang pangkat.
\end{variasi}

% ---------------- SOAL 1.1.6 ----------------
\soalno{1.1.6}{membandingkan dua pangkat}
\begin{soalbox}{}
Bandingkan $2^{10}$ dan $10^3$; mana yang lebih besar?
\end{soalbox}

\jawab
\langkah{1} Cara paling jujur: hitung keduanya.
\[
2^{10} = 1024, \qquad 10^3 = 1000.
\]
\langkah{2} Bandingkan: $1024 > 1000$.
\jadi $2^{10} > 10^3$ (selisihnya hanya $24$, sangat tipis!).

\begin{center}
\begin{tikzpicture}
\draw[fill=birumuda,draw=biru] (0,0) rectangle (10.24,0.6);
\node[right] at (10.24,0.3) {\small $2^{10}=1024$};
\draw[fill=oranyemuda,draw=oranye] (0,-0.9) rectangle (10.0,-0.3);
\node[right] at (10.0,-0.6) {\small $10^3=1000$};
\draw[dashed,gray] (10.0,-1.1) -- (10.0,0.8);
\node[gray,font=\scriptsize] at (10.5,0.95) {beda 24};
\end{tikzpicture}\\
{\small Dua batang hampir sama panjang --- itulah kenapa $2^{10}\approx 10^3$
sering dipakai insinyur ($1$ KB $\approx 1000$ byte).}
\end{center}

\begin{jebakan}
\begin{itemize}[leftmargin=*,topsep=2pt]
  \item \textbf{Menebak dari ``angka lebih besar''.} Pangkat $10$ vs basis $10$:
  tidak bisa ditebak dari sekilas. Basis besar (10) bisa kalah oleh pangkat besar
  (10).
  \item \textbf{Menyamakan basis secara paksa} padahal $2$ dan $10$ tak sepangkat
  bulat. Untuk soal seperti ini, hitung langsung lebih aman.
\end{itemize}
\end{jebakan}

\begin{variasi}
Bandingkan $2^{30}$ dan $10^9$ (petunjuk: $2^{30}=(2^{10})^3\approx(10^3)^3=10^9$,
sedikit lebih besar), atau $3^{20}$ vs $2^{30}$ (samakan lewat $\log$). Ide
umum: \emph{pangkatkan supaya bisa dibandingkan}.
\end{variasi}

% ---------------- SOAL 1.1.7 ----------------
\soalno{1.1.7}{peluruhan — bola memantul}
\begin{soalbox}{}
Sebuah bola dijatuhkan; setiap pantulan mencapai $0{,}8$ kali tinggi
sebelumnya. Berapa tinggi pantulan ke-5 jika tinggi awalnya $2$ m?
\end{soalbox}

\begin{ingat}{Model peluruhan geometri}
Kalau tiap langkah dikalikan faktor tetap $r$ (di sini $r=0{,}8$), maka nilai
setelah $n$ langkah adalah $\text{awal}\times r^{\,n}$. Karena $0<r<1$, nilainya
\emph{menyusut} (peluruhan).
\end{ingat}

\jawab
\langkah{1} Tinggi awal (sebelum pantulan) $=2$ m. Tinggi \emph{pantulan ke-$n$}:
\[
h_n = 2 \times 0{,}8^{\,n}.
\]
\langkah{2} Untuk pantulan ke-5, masukkan $n=5$:
\[
h_5 = 2\times 0{,}8^{5}.
\]
\langkah{3} Hitung $0{,}8^5$ bertahap:
$0{,}8^2=0{,}64$; $0{,}8^3=0{,}512$; $0{,}8^4=0{,}4096$; $0{,}8^5=0{,}32768$.
\langkah{4} $h_5 = 2\times 0{,}32768 = 0{,}65536$ m.
\jadi tinggi pantulan ke-5 $\approx 0{,}66$ m.

\begin{center}
\begin{tikzpicture}[scale=0.9]
\draw[->] (0,0) -- (9,0) node[right]{lintasan};
\draw[->] (0,0) -- (0,4.2) node[above]{tinggi (m)};
\foreach \x/\h/\lab in {0/4/2{,}0, 2/3.2/1{,}6, 3.6/2.56/1{,}28, 4.9/2.05/1{,}02, 6/1.64/0{,}82, 6.9/1.31/0{,}66}{
  \node[circle,fill=biru,inner sep=1pt] at (\x,\h) {};
}
\draw[biru,thick] (0,4) .. controls (1,0) and (1,0) .. (2,3.2)
  .. controls (2.8,0) and (2.8,0) .. (3.6,2.56)
  .. controls (4.25,0) and (4.25,0) .. (4.9,2.05)
  .. controls (5.45,0) and (5.45,0) .. (6,1.64)
  .. controls (6.45,0) and (6.45,0) .. (6.9,1.31);
\node[font=\scriptsize] at (0,4.4) {$2{,}0$};
\node[font=\scriptsize,oranye] at (6.9,1.6) {ke-5};
\end{tikzpicture}\\
{\small Tiap puncak $=0{,}8\times$ puncak sebelumnya. Pantulan mengecil terus
tapi \emph{tidak pernah nol}.}
\end{center}

\begin{jebakan}
\begin{itemize}[leftmargin=*,topsep=2pt]
  \item \textbf{Salah menghitung indeks.} ``Pantulan ke-5'' berarti pangkat $5$.
  Kalau soal bilang ``tinggi setelah menyentuh lantai ke-5'', maknanya juga
  pangkat 5 --- tapi baca hati-hati, kadang yang ditanya tinggi \emph{awal}
  (pangkat 0).
  \item \textbf{Mengurangi $0{,}8$ tiap kali} ($2-0{,}8-0{,}8\dots$). Ini
  \emph{perkalian} berulang, bukan pengurangan.
  \item \textbf{Membulatkan terlalu awal.} $0{,}8^2=0{,}64$ eksak; jangan
  dibulatkan ke $0{,}6$ di tengah jalan.
\end{itemize}
\end{jebakan}

\begin{variasi}
\begin{itemize}[leftmargin=*,topsep=2pt]
  \item \textbf{Total jarak tempuh} bola sampai berhenti $\to$ deret geometri tak
  hingga (dibahas di Bab 2).
  \item \textbf{Peluruhan lain:} obat berkurang $10\%$ tiap jam ($r=0{,}9$), zat
  radioaktif tinggal setengah tiap periode ($r=0{,}5$), harga mobil turun $15\%$
  per tahun.
\end{itemize}
\end{variasi}

% ---------------- SOAL 1.1.8 ----------------
\soalno{1.1.8}{akar sebagai pangkat pecahan}
\begin{soalbox}{}
Sederhanakan $\sqrt{a}\cdot\sqrt[3]{a}$ dalam bentuk pangkat tunggal.
\end{soalbox}

\jawab
\langkah{1} Ubah semua akar menjadi pangkat pecahan (sifat 8):
$\sqrt{a}=a^{1/2}$ dan $\sqrt[3]{a}=a^{1/3}$.
\langkah{2} Kalikan pangkat sama-basis $\Rightarrow$ pangkat \emph{dijumlah}
(sifat 1). Samakan penyebut:
\[
a^{1/2}\cdot a^{1/3} = a^{\frac12+\frac13} = a^{\frac{3}{6}+\frac{2}{6}} = a^{5/6}.
\]
\jadi $\sqrt a\cdot\sqrt[3]a = a^{5/6}$ (atau $\sqrt[6]{a^5}$).

\begin{jebakan}
\begin{itemize}[leftmargin=*,topsep=2pt]
  \item \textbf{Mengalikan penyebut akar} ($\sqrt{}\cdot\sqrt[3]{}=\sqrt[6]{}$
  seolah $2\times3=6$ pada indeks). Yang benar: \emph{jumlahkan pangkat}
  $\frac12+\frac13$, bukan kalikan indeks.
  \item \textbf{Lupa menyamakan penyebut} sebelum menjumlah pecahan.
\end{itemize}
\end{jebakan}

\begin{variasi}
$\dfrac{\sqrt a}{\sqrt[3]{a}} = a^{1/2-1/3}=a^{1/6}$; atau
$\sqrt{a\sqrt a}=\sqrt{a\cdot a^{1/2}}=\sqrt{a^{3/2}}=a^{3/4}$ (akar bertingkat).
\end{variasi}

% ---------------- SOAL 1.1.9 ----------------
\soalno{1.1.9}{pangkat negatif pada pecahan}
\begin{soalbox}{}
Tentukan nilai $\left(\tfrac{1}{2}\right)^{-3}$.
\end{soalbox}

\jawab
\langkah{1} Pangkat negatif $=$ balik pecahannya (sifat 7 pada pecahan):
$\left(\frac ab\right)^{-n}=\left(\frac ba\right)^{n}$.
\[
\left(\tfrac12\right)^{-3} = \left(\tfrac21\right)^{3} = 2^3 = 8.
\]
\jadi $\left(\tfrac12\right)^{-3} = 8$.

\begin{jebakan}
\begin{itemize}[leftmargin=*,topsep=2pt]
  \item \textbf{Mengira hasilnya pecahan/negatif}, mis.\ $-\frac18$ atau
  $\frac18$. Pangkat negatif pada pecahan kecil justru menghasilkan bilangan
  \emph{besar}, karena membalik pecahan.
  \item \textbf{Membalik tapi lupa pangkat}: $\left(\frac12\right)^{-3}\neq 2$.
  Setelah dibalik jadi $2$, masih harus dipangkatkan $3$.
\end{itemize}
\end{jebakan}

\begin{variasi}
$\left(\tfrac23\right)^{-2}=\left(\tfrac32\right)^2=\tfrac94$;
$\left(0{,}5\right)^{-4}=2^4=16$ (ubah desimal ke pecahan dulu).
\end{variasi}

% ---------------- SOAL 1.1.10 ----------------
\soalno{1.1.10}{model pertumbuhan tabungan}
\begin{soalbox}{}
Sebuah tabungan menjadi $1{,}05$ kali lipat tiap tahun. Tuliskan saldo setelah
$t$ tahun jika saldo awal Rp$1.000.000$.
\end{soalbox}

\begin{ingat}{Rumus umum pertumbuhan}
Nilai($t$) $=$ Nilai awal $\times (\text{faktor pengali})^{t}$. Bila naik $p\%$
per periode, faktor pengalinya $=1+\frac{p}{100}$. Naik $5\%\Rightarrow \times1{,}05$.
\end{ingat}

\jawab
\langkah{1} Kenali komponennya: nilai awal $M_0 = 1.000.000$; faktor pengali
$=1{,}05$ (karena naik $5\%$); pangkat = banyak tahun $t$.
\langkah{2} Susun modelnya:
\[
\boxed{\,M(t) = 1.000.000 \times 1{,}05^{\,t}\,}
\]
\jadi saldo setelah $t$ tahun adalah $1.000.000\times 1{,}05^{t}$ rupiah.
(Contoh: $t=10\Rightarrow 1.000.000\times1{,}05^{10}\approx \text{Rp}\,1.628.895$.)

\begin{jebakan}
\begin{itemize}[leftmargin=*,topsep=2pt]
  \item \textbf{Menulis faktor $0{,}05$} bukan $1{,}05$. Faktor $0{,}05$ berarti
  saldo \emph{menyusut} tinggal $5\%$ --- keliru. Naik $5\%$ berarti menjadi
  $100\%+5\%=105\%=1{,}05$ kali.
  \item \textbf{Model linear.} Menulis $1.000.000 + 50.000t$ (bunga tunggal).
  Soal ini bunga \emph{majemuk} (dilipatkan), jadi pangkat, bukan tambah tetap.
  Bandingkan di Bab 9.
  \item \textbf{Meletakkan $t$ sebagai basis} ($1{,}05\cdot t$). $t$ ada di
  \emph{pangkat}, bukan dikalikan.
\end{itemize}
\end{jebakan}

\begin{variasi}
Turun/peluruhan: ``nilai mesin menyusut $8\%$/tahun'' $\to M_0\times 0{,}92^t$.
Ditambah pertanyaan waktu: ``kapan saldo jadi 2 juta?'' $\to$ butuh logaritma
(Subbab 1.3). Ini persis kerangka \emph{studi kasus} bab ini.
\end{variasi}

% ############################################################
\section*{Subbab 1.2 — Persamaan dan Pertidaksamaan Eksponen}
\addcontentsline{toc}{section}{Subbab 1.2 — Persamaan \& Pertidaksamaan Eksponen}
% ############################################################

\begin{ingat}{Prinsip ``samakan basis''}
Jika $a^{f(x)} = a^{g(x)}$ dengan $a>0,\ a\neq1$, maka $f(x)=g(x)$.
Untuk pertidaksamaan:
\begin{itemize}[leftmargin=*,topsep=1pt]
  \item basis $a>1$: tanda \textbf{tetap} ($a^{f}>a^{g}\Rightarrow f>g$);
  \item basis $0<a<1$: tanda \textbf{dibalik} ($a^{f}>a^{g}\Rightarrow f<g$).
\end{itemize}
Untuk bentuk kuadrat (mis.\ $9^x-4\cdot3^x+3$), pakai \emph{pemisalan} $p=a^x$.
\end{ingat}

% ---------------- SOAL 1.2.1 ----------------
\soalno{1.2.1}{persamaan eksponen dasar}
\begin{soalbox}{}
Selesaikan $3^{2x+1} = 81$.
\end{soalbox}

\jawab
\langkah{1} Samakan basis. Ubah ruas kanan ke basis $3$: $81 = 3^4$.
\langkah{2} Persamaan menjadi $3^{2x+1} = 3^{4}$. Karena basis sama, pangkat harus
sama:
\[
2x+1 = 4.
\]
\langkah{3} Selesaikan: $2x=3\Rightarrow x=\tfrac32$.
\jadi $x=\dfrac32$.

\begin{jebakan}
\begin{itemize}[leftmargin=*,topsep=2pt]
  \item \textbf{Salah menyatakan $81$.} $81=3^4$ (bukan $3^3=27$ atau $9^2$ yang
  lupa disamakan basisnya). Hafalkan pangkat kecil: $3^1,3^2,3^3,3^4=3,9,27,81$.
  \item \textbf{Menyamakan pangkat sebelum basis sama.} Tidak boleh langsung
  ``$2x+1=81$''. Ubah dulu $81$ ke basis $3$.
\end{itemize}
\end{jebakan}

\begin{variasi}
$3^{2x+1}=\tfrac1{27}$ ($=3^{-3}$, jawaban negatif); $3^{2x+1}=1$ ($=3^0$).
Selalu ubah ruas kanan ke pangkat basis yang sama --- termasuk $1=a^0$.
\end{variasi}

% ---------------- SOAL 1.2.2 ----------------
\soalno{1.2.2}{persamaan kuadrat di pangkat}
\begin{soalbox}{}
Selesaikan $5^{x^2-3x} = 5^{-2}$.
\end{soalbox}

\jawab
\langkah{1} Basis sudah sama ($5$), langsung samakan pangkat:
\[
x^2 - 3x = -2.
\]
\langkah{2} Pindahkan semua ke satu ruas: $x^2-3x+2=0$.
\langkah{3} Faktorkan: $(x-1)(x-2)=0 \Rightarrow x=1$ atau $x=2$.
\jadi $x=1$ atau $x=2$.

\begin{jebakan}
\begin{itemize}[leftmargin=*,topsep=2pt]
  \item \textbf{Lupa memindah $-2$.} Menulis $x^2-3x=-2$ lalu memfaktorkan
  langsung tanpa membuat ruas kanan nol. Persamaan kuadrat harus $=0$ dulu.
  \item \textbf{Hanya mengambil satu akar.} Persamaan kuadrat bisa punya dua
  solusi; keduanya sah karena tidak melanggar syarat (pangkat boleh bilangan
  apa saja).
\end{itemize}
\end{jebakan}

\begin{variasi}
$2^{x^2-1}=8$ (ubah $8=2^3$), $7^{x^2}=7^{2x+3}$. Pola: samakan basis $\to$
persamaan kuadrat $\to$ faktorkan.
\end{variasi}

% ---------------- SOAL 1.2.3 ----------------
\soalno{1.2.3}{dua basis, satu ``leluhur'' bersama}
\begin{soalbox}{}
Selesaikan $4^{x} = 8^{x-1}$.
\end{soalbox}

\begin{ingat}{Cari basis prima bersama}
Kalau dua basis berbeda ($4$ dan $8$) tapi punya ``leluhur'' prima yang sama
(di sini $2$), ubah keduanya ke basis prima itu: $4=2^2,\ 8=2^3$.
\end{ingat}

\jawab
\langkah{1} Tulis ke basis $2$: $4^x = (2^2)^x = 2^{2x}$; $8^{x-1}=(2^3)^{x-1}=2^{3(x-1)}$.
\langkah{2} Samakan pangkat: $2x = 3(x-1)$.
\langkah{3} Selesaikan: $2x = 3x - 3 \Rightarrow -x = -3 \Rightarrow x = 3$.
\jadi $x=3$. (Cek: $4^3=64$ dan $8^{2}=64$. \checkmark)

\begin{jebakan}
\begin{itemize}[leftmargin=*,topsep=2pt]
  \item \textbf{Lupa mengalikan seluruh $(x-1)$.} $(2^3)^{x-1}=2^{3x-3}$, bukan
  $2^{3x-1}$. Pangkat $3$ mengali \emph{semua} isi kurung.
  \item \textbf{Menyamakan pangkat sebelum basis sama} (``$x=x-1$'').
\end{itemize}
\end{jebakan}

\begin{variasi}
$27^{x}=9^{x+1}$ (basis $3$), $\left(\tfrac14\right)^x = 8^{x+2}$ (campur pangkat
positif--negatif dari $2$). Selalu turunkan ke basis prima bersama.
\end{variasi}

% ---------------- SOAL 1.2.4 ----------------
\soalno{1.2.4}{pemisalan (persamaan kuadrat tersembunyi)}
\begin{soalbox}{}
Selesaikan $9^{x} - 4\cdot 3^{x} + 3 = 0$.
\end{soalbox}

\begin{ingat}{Kapan memakai pemisalan}
Jika muncul $a^{2x}$ (atau $a^{2}{}^{x}$) \emph{dan} $a^{x}$ bersama, misalkan
$p=a^{x}$. Karena $9^x=(3^2)^x=(3^x)^2=p^2$, bentuknya berubah jadi kuadrat biasa
dalam $p$.
\end{ingat}

\jawab
\langkah{1} Nyatakan $9^x$ lewat $3^x$: $9^x=(3^x)^2$. Misalkan $p=3^x$
(catat: $p>0$ selalu).
\langkah{2} Persamaan menjadi $p^2 - 4p + 3 = 0$.
\langkah{3} Faktorkan: $(p-1)(p-3)=0 \Rightarrow p=1$ atau $p=3$.
\langkah{4} Kembalikan ke $x$:
\[
3^x = 1 = 3^0 \Rightarrow x=0; \qquad 3^x = 3 = 3^1 \Rightarrow x=1.
\]
\jadi $x=0$ atau $x=1$.

\begin{jebakan}
\begin{itemize}[leftmargin=*,topsep=2pt]
  \item \textbf{Lupa mengembalikan ke $x$.} Menjawab ``$p=1,3$'' bukan jawaban
  akhir --- yang ditanya $x$.
  \item \textbf{Membuang solusi valid.} Di sini $p=1$ sah karena $3^x=1$ punya
  jawaban $x=0$. (Solusi hanya dibuang jika $p\le0$, sebab $3^x$ tak pernah
  negatif atau nol.)
  \item \textbf{Salah mengurai $9^x$.} $9^x\neq 3\cdot3^x$ dan $\neq 3^{x}\cdot 3$.
  Yang benar $9^x=(3^x)^2$.
\end{itemize}
\end{jebakan}

\begin{variasi}
$4^x-6\cdot2^x+8=0$, $2^{2x}-5\cdot2^x+4=0$, atau bentuk dengan solusi
terbuang: $4^x-2^x-2=0\Rightarrow p^2-p-2=0\Rightarrow(p-2)(p+1)=0$; $p=-1$
\emph{dibuang} (mustahil), sisa $p=2\Rightarrow x=1$.
\end{variasi}

% ---------------- SOAL 1.2.5 ----------------
\soalno{1.2.5}{pertidaksamaan, basis > 1}
\begin{soalbox}{}
Tentukan himpunan penyelesaian $2^{x} \le 16$.
\end{soalbox}

\jawab
\langkah{1} Samakan basis: $16=2^4$, sehingga $2^x \le 2^4$.
\langkah{2} Basis $2>1$ $\Rightarrow$ tanda \textbf{dipertahankan}:
\[
x \le 4.
\]
\jadi HP $=\{x \mid x\le 4\}$.

\begin{jebakan}
\begin{itemize}[leftmargin=*,topsep=2pt]
  \item \textbf{Membalik tanda tanpa alasan.} Tanda hanya dibalik kalau basis
  $<1$. Di sini basis $2>1$, tanda tetap.
\end{itemize}
\end{jebakan}

\begin{variasi}
$2^x < 32$, $2^{3x-1}\ge 8$, atau $2^{x^2}\le 16$ (jadi $x^2\le4\Rightarrow
-2\le x\le2$). Perhatikan kapan muncul pertidaksamaan kuadrat.
\end{variasi}

% ---------------- SOAL 1.2.6 ----------------
\soalno{1.2.6}{pertidaksamaan, basis < 1 (tanda dibalik)}
\begin{soalbox}{}
Tentukan himpunan penyelesaian $\left(\tfrac12\right)^{2x-1} \ge 8$.
\end{soalbox}

\begin{ingat}{Kenapa tanda dibalik saat basis $<1$?}
Fungsi $\left(\tfrac12\right)^x$ \emph{turun}: makin besar $x$, makin
\emph{kecil} nilainya. Jadi ``nilai lebih besar'' justru datang dari ``pangkat
lebih kecil'' --- itulah sebabnya arah pertidaksamaan terbalik.
\end{ingat}

\jawab
\langkah{1} Samakan basis ke $\tfrac12$: $8 = 2^3 = \left(\tfrac12\right)^{-3}$.
Jadi $\left(\tfrac12\right)^{2x-1} \ge \left(\tfrac12\right)^{-3}$.
\langkah{2} Basis $\tfrac12<1$ $\Rightarrow$ tanda \textbf{dibalik}:
\[
2x-1 \le -3.
\]
\langkah{3} Selesaikan: $2x \le -2 \Rightarrow x \le -1$.
\jadi HP $=\{x\mid x\le -1\}$.

\begin{jebakan}
\begin{itemize}[leftmargin=*,topsep=2pt]
  \item \textbf{Lupa membalik tanda} --- kesalahan nomor satu di topik ini.
  Sebelum menyamakan pangkat, \emph{selalu cek besar basis}.
  \item \textbf{Salah menulis $8$ sebagai pangkat $\tfrac12$.}
  $8=\left(\tfrac12\right)^{-3}$ (pangkat negatif membalik $\tfrac12$ jadi $2$,
  lalu $2^3=8$).
\end{itemize}
\end{jebakan}

\begin{variasi}
$\left(\tfrac13\right)^x>\tfrac19$ (jawaban $x<2$), $0{,}25^{x}\le 0{,}5$
(ubah ke basis $\tfrac12$). Ciri khas: basis pecahan $\Rightarrow$ waspada balik
tanda.
\end{variasi}

% ---------------- SOAL 1.2.7 ----------------
\soalno{1.2.7}{basis sama sejak awal}
\begin{soalbox}{}
Selesaikan $7^{x+2} = 7^{3x-4}$.
\end{soalbox}

\jawab
\langkah{1} Basis sudah sama ($7$), langsung samakan pangkat:
$x+2 = 3x-4$.
\langkah{2} Kumpulkan $x$: $2+4 = 3x - x \Rightarrow 6 = 2x \Rightarrow x=3$.
\jadi $x=3$.

\begin{jebakan}
\begin{itemize}[leftmargin=*,topsep=2pt]
  \item \textbf{Salah pindah ruas} (kesalahan tanda saat memindahkan $x$ dan
  angka). Rapikan: variabel ke kanan, konstanta ke kiri.
\end{itemize}
\end{jebakan}

\begin{variasi}
Ditambah pemfaktoran: $7^{x+2}=49^{x-1}=7^{2x-2}$; atau ruas kanan $1$:
$7^{x+2}=1=7^0\Rightarrow x=-2$.
\end{variasi}

% ---------------- SOAL 1.2.8 ----------------
\soalno{1.2.8}{persamaan simetri $a^x+a^{-x}$}
\begin{soalbox}{}
Jika $2^{x}+2^{-x} = \tfrac52$, tentukan nilai $2^{x}$.
\end{soalbox}

\begin{ingat}{Trik pemisalan pada bentuk simetri}
$2^{-x}=\dfrac{1}{2^{x}}$. Misalkan $p=2^x$ ($p>0$); maka $2^{-x}=\frac1p$.
Persamaan berubah menjadi $p+\frac1p=\frac52$, yang bisa dijadikan kuadrat dengan
mengalikan $p$.
\end{ingat}

\jawab
\langkah{1} Misalkan $p=2^x$, sehingga $2^{-x}=\tfrac1p$. Persamaan:
$p+\tfrac1p = \tfrac52$.
\langkah{2} Kalikan kedua ruas dengan $2p$ untuk menghapus pecahan:
\[
2p^2 + 2 = 5p \Rightarrow 2p^2 - 5p + 2 = 0.
\]
\langkah{3} Faktorkan: $(2p-1)(p-2)=0 \Rightarrow p=\tfrac12$ atau $p=2$.
\langkah{4} Karena $p=2^x$: $2^x=\tfrac12$ (yakni $x=-1$) atau $2^x=2$ (yakni
$x=1$).
\jadi $2^x = 2$ atau $2^x=\dfrac12$.

\begin{jebakan}
\begin{itemize}[leftmargin=*,topsep=2pt]
  \item \textbf{Membuang salah satu solusi tanpa alasan.} Keduanya positif, jadi
  \emph{keduanya sah}. (Kedua nilai $x=1$ dan $x=-1$ memang saling ``pasangan''
  karena bentuknya simetris.)
  \item \textbf{Salah kali silang.} Saat mengalikan $p$, ingat $\frac52\cdot p$
  dan $\frac1p\cdot p=1$. Rapikan sebelum memfaktorkan.
\end{itemize}
\end{jebakan}

\begin{variasi}
Diberi $2^x+2^{-x}=3$, ditanya $4^x+4^{-x}$ (kuadratkan: $=9-2=7$). Atau
$3^x+3^{-x}=\frac{10}{3}$. Pola simetri ini favorit soal SNBT.
\end{variasi}

% ---------------- SOAL 1.2.9 ----------------
\soalno{1.2.9}{menyusun ulang sebelum menyamakan}
\begin{soalbox}{}
Selesaikan $25^{x} = 5\cdot 5^{x}$.
\end{soalbox}

\jawab
\langkah{1} Ubah semua ke basis $5$: $25^x=(5^2)^x=5^{2x}$; ruas kanan
$5\cdot5^x = 5^{1}\cdot5^{x} = 5^{x+1}$ (sifat 1).
\langkah{2} Persamaan: $5^{2x} = 5^{x+1}$. Samakan pangkat:
\[
2x = x+1 \Rightarrow x = 1.
\]
\jadi $x=1$. (Cek: $25^1=25$ dan $5\cdot5^1=25$. \checkmark)

\begin{jebakan}
\begin{itemize}[leftmargin=*,topsep=2pt]
  \item \textbf{Menganggap $5\cdot5^x = 25^x$.} Tidak! $5\cdot5^x=5^{x+1}$
  (pangkat \emph{ditambah} 1), sedangkan $25^x=5^{2x}$.
  \item \textbf{Lupa menulis $5=5^1$} sehingga sifat perkalian pangkat tak
  terlihat.
\end{itemize}
\end{jebakan}

\begin{variasi}
$4^x = 2\cdot 8^{x}$, $9\cdot3^x = 27^{x}$. Rahasianya sama: rapikan tiap ruas
menjadi \emph{satu} pangkat basis prima, baru samakan.
\end{variasi}

% ---------------- SOAL 1.2.10 ----------------
\soalno{1.2.10}{menyusun persamaan dari studi kasus}
\begin{soalbox}{}
Berapa lama (dalam bulan) subscriber pada studi kasus menembus $100.000$?
(Cukup nyatakan dalam bentuk \emph{persamaan eksponen}-nya.)
\end{soalbox}

\jawab
\langkah{1} Ingat model studi kasus: $U_n = 1000\times 1{,}5^{n}$.
\langkah{2} Syarat ``menembus $100.000$'':
\[
1000\times 1{,}5^{\,n} = 100.000.
\]
\langkah{3} Bagi kedua ruas dengan $1000$ agar rapi:
\[
\boxed{\,1{,}5^{\,n} = 100\,}
\]
Inilah persamaan eksponen yang diminta. \emph{Menyelesaikannya untuk $n$
memerlukan logaritma} (Subbab 1.3), dan hasilnya $n=\,{}^{1{,}5}\!\log100\approx
11{,}4$, sehingga tercapai pada bulan ke-12.

\begin{jebakan}
\begin{itemize}[leftmargin=*,topsep=2pt]
  \item \textbf{Langsung membagi $100000/1000=100$ lalu $100/1{,}5$.} Salah ---
  $n$ ada di \emph{pangkat}, tak bisa dibagi begitu saja.
  \item \textbf{Lupa mencantumkan nilai awal $1000$.} Model harus berangkat dari
  $1000$, bukan langsung $1{,}5^n=100000$.
\end{itemize}
\end{jebakan}

\begin{variasi}
``Kapan menembus $1$ juta?'' $\Rightarrow 1{,}5^n=1000$; ``kapan menjadi $2$ kali
lipat?'' $\Rightarrow 1{,}5^n=2$. Semua menyisakan $n=\,{}^{1{,}5}\!\log(\cdots)$.
\end{variasi}

% ############################################################
\section*{Subbab 1.3 — Logaritma}
\addcontentsline{toc}{section}{Subbab 1.3 — Logaritma}
% ############################################################

\begin{ingat}{Definisi \& sifat logaritma}
Logaritma adalah \emph{invers} eksponen:
\[
{}^{a}\!\log x = y \iff a^{y}=x \quad (a>0,\ a\neq1,\ x>0).
\]
Baca ${}^{a}\!\log x$ sebagai: ``$a$ dipangkatkan \emph{berapa} agar jadi $x$?''
Sifat-sifatnya:
\[
\begin{array}{ll}
(1)\ {}^{a}\!\log(xy)={}^{a}\!\log x+{}^{a}\!\log y &
(2)\ {}^{a}\!\log\dfrac xy={}^{a}\!\log x-{}^{a}\!\log y\\[7pt]
(3)\ {}^{a}\!\log x^n = n\,{}^{a}\!\log x &
(4)\ {}^{a}\!\log a=1,\ \ {}^{a}\!\log1=0\\[7pt]
(5)\ {}^{a}\!\log x=\dfrac{{}^{c}\!\log x}{{}^{c}\!\log a}\ (\text{ganti basis}) &
(6)\ a^{{}^{a}\!\log x}=x
\end{array}
\]
\textbf{Inti:} logaritma mengubah \emph{kali}$\to$\emph{tambah},
\emph{bagi}$\to$\emph{kurang}, \emph{pangkat}$\to$\emph{kali}.
\end{ingat}

% ---------------- SOAL 1.3.1 ----------------
\soalno{1.3.1}{menghitung nilai logaritma}
\begin{soalbox}{}
Hitung ${}^{2}\!\log 16 + {}^{4}\!\log 64$.
\end{soalbox}

\jawab
\langkah{1} Kerjakan tiap logaritma \emph{terpisah}, dengan menuliskan argumen
sebagai pangkat basisnya.
\langkah{2} ${}^{2}\!\log 16$: ``$2$ pangkat berapa $=16$?'' Karena $16=2^4$,
maka ${}^{2}\!\log16=4$.
\langkah{3} ${}^{4}\!\log 64$: ``$4$ pangkat berapa $=64$?'' Karena $64=4^3$,
maka ${}^{4}\!\log64=3$.
\langkah{4} Jumlahkan: $4+3=7$.
\jadi ${}^{2}\!\log16+{}^{4}\!\log64 = 7$.

\begin{jebakan}
\begin{itemize}[leftmargin=*,topsep=2pt]
  \item \textbf{Menyamakan basis paksa yang salah.} Basis berbeda ($2$ dan $4$)
  tidak boleh dijumlah begitu saja seolah satu. Hitung sendiri-sendiri.
  \item \textbf{Salah kenali pangkat.} $64$ bisa $2^6, 4^3, 8^2$ --- pilih yang
  \emph{sesuai basisnya}. Untuk ${}^{4}\!\log$, tulis $64=4^3$.
\end{itemize}
\end{jebakan}

\begin{variasi}
${}^{3}\!\log 27 + {}^{9}\!\log 81$, ${}^{5}\!\log125 - {}^{2}\!\log32$. Selalu
nyatakan argumen sebagai pangkat basisnya.
\end{variasi}

% ---------------- SOAL 1.3.2 ----------------
\soalno{1.3.2}{menyatakan log dalam variabel}
\begin{soalbox}{}
Jika ${}^{2}\!\log 3 = p$, nyatakan ${}^{2}\!\log 18$ dalam $p$.
\end{soalbox}

\begin{ingat}{Strategi ``faktorkan argumennya''}
Kalau diminta menyatakan logaritma dalam variabel, \emph{faktorkan angka di
dalam log} menjadi campuran (a) bilangan yang sama dengan basis, dan (b) bilangan
yang sudah punya variabel.
\end{ingat}

\jawab
\langkah{1} Faktorkan $18$: $18 = 2\times 9 = 2\times 3^2$.
\langkah{2} Pakai sifat (1) lalu (3):
\[
{}^{2}\!\log 18 = {}^{2}\!\log(2\cdot3^2) = {}^{2}\!\log 2 + {}^{2}\!\log 3^2
= {}^{2}\!\log 2 + 2\,{}^{2}\!\log 3.
\]
\langkah{3} Substitusi ${}^{2}\!\log2=1$ (sifat 4) dan ${}^{2}\!\log3=p$:
\[
= 1 + 2p.
\]
\jadi ${}^{2}\!\log 18 = 1+2p$.

\begin{jebakan}
\begin{itemize}[leftmargin=*,topsep=2pt]
  \item \textbf{Lupa ${}^{2}\!\log2=1$.} Sering ditinggal sebagai
  ``${}^{2}\!\log2$'' padahal nilainya $1$.
  \item \textbf{Salah faktor.} $18=2\cdot3^2$, bukan $2\cdot9$ yang dibiarkan
  ($9$ harus jadi $3^2$ agar muncul $p$).
  \item \textbf{Menganggap ${}^{2}\!\log18={}^{2}\!\log2\cdot{}^{2}\!\log9$.}
  Log \emph{perkalian} jadi \emph{penjumlahan}, bukan perkalian log.
\end{itemize}
\end{jebakan}

\begin{variasi}
``${}^{2}\!\log3=p$, nyatakan ${}^{2}\!\log12$'' ($=2+p$), atau dua variabel:
``${}^{2}\!\log3=p,\ {}^{2}\!\log5=q$, nyatakan ${}^{2}\!\log 45$''
($45=3^2\cdot5\Rightarrow 2p+q$).
\end{variasi}

% ---------------- SOAL 1.3.3 ----------------
\soalno{1.3.3}{sifat selisih logaritma}
\begin{soalbox}{}
Hitung ${}^{3}\!\log 81 - {}^{3}\!\log 9$.
\end{soalbox}

\jawab
\textbf{Cara A (hitung terpisah).} ${}^{3}\!\log81=4$ (sebab $3^4=81$);
${}^{3}\!\log9=2$ (sebab $3^2=9$). Selisih $4-2=2$.

\textbf{Cara B (sifat selisih, sifat 2).}
${}^{3}\!\log81 - {}^{3}\!\log9 = {}^{3}\!\log\dfrac{81}{9} = {}^{3}\!\log9 = 2.$
\jadi hasilnya $2$.

\begin{jebakan}
\begin{itemize}[leftmargin=*,topsep=2pt]
  \item \textbf{Membagi nilainya} ($\frac{4}{2}=2$ secara kebetulan benar di
  sini, tapi \emph{itu logika salah}). Selisih \emph{log} $=$ log \emph{bagi};
  bukan log dibagi log.
  \item \textbf{Menganggap ${}^3\!\log81 - {}^3\!\log9={}^3\!\log(81-9)$.} Salah
  --- yang berlaku adalah log \emph{bagi}, bukan log selisih.
\end{itemize}
\end{jebakan}

\begin{variasi}
${}^{2}\!\log40-{}^{2}\!\log5$ ($={}^{2}\!\log8=3$),
${}^{5}\!\log100-{}^{5}\!\log4$ ($={}^{5}\!\log25=2$).
\end{variasi}

% ---------------- SOAL 1.3.4 ----------------
\soalno{1.3.4}{gabungan jumlah \& selisih}
\begin{soalbox}{}
Sederhanakan ${}^{2}\!\log 6 + {}^{2}\!\log 8 - {}^{2}\!\log 3$.
\end{soalbox}

\jawab
\langkah{1} Karena basisnya sama semua ($2$), gabungkan jadi satu log: tambah
$\to$ kali, kurang $\to$ bagi.
\[
{}^{2}\!\log\frac{6\times 8}{3} = {}^{2}\!\log\frac{48}{3} = {}^{2}\!\log 16.
\]
\langkah{2} ${}^{2}\!\log16 = 4$ (sebab $2^4=16$).
\jadi hasilnya $4$.

\begin{jebakan}
\begin{itemize}[leftmargin=*,topsep=2pt]
  \item \textbf{Salah tempatkan yang dibagi.} Yang \emph{dikurang}
  (${}^{2}\!\log3$) masuk ke \emph{penyebut}. Suku bertanda $+$ ke pembilang.
  \item \textbf{Menghitung sebagian, lupa menyederhanakan sisa.} Setelah dapat
  ${}^{2}\!\log16$, jangan berhenti --- nilainya $4$.
\end{itemize}
\end{jebakan}

\begin{variasi}
${}^{3}\!\log 18 + {}^{3}\!\log 6 - {}^{3}\!\log 4$
($=\,{}^{3}\!\log 27=3$), atau tambah pangkat:
${}^{2}\!\log 9 + 2\,{}^{2}\!\log 2 - {}^{2}\!\log\tfrac94$.
\end{variasi}

% ---------------- SOAL 1.3.5 ----------------
\soalno{1.3.5}{persamaan logaritma sederhana}
\begin{soalbox}{}
Tentukan nilai $x$ dari ${}^{2}\!\log x = 5$.
\end{soalbox}

\jawab
\langkah{1} Terjemahkan ke bentuk eksponen memakai definisi
${}^{a}\!\log x=y\iff a^y=x$:
\[
{}^{2}\!\log x = 5 \iff x = 2^{5}.
\]
\langkah{2} $2^5=32$.
\jadi $x=32$.

\begin{jebakan}
\begin{itemize}[leftmargin=*,topsep=2pt]
  \item \textbf{Membalik posisi} ($x=5^2=25$). Ingat pola: basis (bawah) tetap
  jadi basis pangkat; nilai log jadi \emph{pangkatnya}. Basis $2$, pangkat $5$
  $\Rightarrow 2^5$.
  \item \textbf{Lupa syarat $x>0$.} Untuk persamaan log, hasil harus positif ---
  di sini $32>0$, aman.
\end{itemize}
\end{jebakan}

\begin{variasi}
${}^{3}\!\log x=4\Rightarrow x=81$; ${}^{2}\!\log x=-3\Rightarrow x=2^{-3}=\tfrac18$
(pangkat negatif $\to$ hasil pecahan, tetap positif).
\end{variasi}

% ---------------- SOAL 1.3.6 ----------------
\soalno{1.3.6}{sifat rantai (ganti basis)}
\begin{soalbox}{}
Hitung ${}^{5}\!\log 10 \cdot {}^{10}\!\log 25$ (gunakan sifat rantai).
\end{soalbox}

\begin{ingat}{Sifat rantai logaritma}
${}^{a}\!\log b \cdot {}^{b}\!\log c = {}^{a}\!\log c$. Basis-tengah yang sama
($b$) ``saling meniadakan'', menyisakan basis awal dan argumen akhir. Ini
turunan dari sifat ganti basis (5).
\end{ingat}

\jawab
\langkah{1} Perhatikan basis-tengah: argumen log pertama ($10$) = basis log kedua
($10$). Rantai bekerja:
\[
{}^{5}\!\log 10 \cdot {}^{10}\!\log 25 = {}^{5}\!\log 25.
\]
\langkah{2} ${}^{5}\!\log25 = {}^{5}\!\log5^2 = 2$.
\jadi hasilnya $2$.

\begin{jebakan}
\begin{itemize}[leftmargin=*,topsep=2pt]
  \item \textbf{Mengalikan nilai desimalnya} (${}^5\!\log10\approx1{,}43$ kali
  ${}^{10}\!\log25\approx1{,}40$) --- boros dan rawan pembulatan. Sifat rantai
  memberi hasil eksak $2$.
  \item \textbf{Salah urutan rantai.} Rantai butuh pola
  ${}^{a}\!\log b\cdot{}^{b}\!\log c$; pastikan argumen pertama = basis kedua.
\end{itemize}
\end{jebakan}

\begin{variasi}
${}^{2}\!\log3\cdot{}^{3}\!\log8$ ($={}^{2}\!\log8=3$), atau rantai tiga suku:
${}^{a}\!\log b\cdot{}^{b}\!\log c\cdot{}^{c}\!\log a = 1$.
\end{variasi}

% ---------------- SOAL 1.3.7 ----------------
\soalno{1.3.7}{logaritma basis 10 (log biasa)}
\begin{soalbox}{}
Diketahui $\log 2 = 0{,}301$. Hitung $\log 5$.
\end{soalbox}

\begin{ingat}{Trik $5 = \tfrac{10}{2}$}
Untuk log basis $10$ (ditulis $\log$ tanpa basis), $\log10=1$. Karena
$5=\frac{10}{2}$, gunakan sifat selisih. Trik ini menghindari kebutuhan tabel.
\end{ingat}

\jawab
\langkah{1} Tulis $5=\dfrac{10}{2}$.
\langkah{2} Sifat selisih (2): $\log5 = \log\dfrac{10}{2} = \log10 - \log2$.
\langkah{3} $\log10=1$ dan $\log2=0{,}301$:
\[
\log5 = 1 - 0{,}301 = 0{,}699.
\]
\jadi $\log5 = 0{,}699$.

\begin{jebakan}
\begin{itemize}[leftmargin=*,topsep=2pt]
  \item \textbf{Mengira $\log5=\tfrac{\log10}{\log2}$.} Itu \emph{ganti basis},
  bukan yang diminta. $5=10\div2$ berarti log-nya \emph{dikurang}, bukan dibagi.
  \item \textbf{Lupa $\log10=1$.} Ini fakta wajib untuk log basis $10$.
\end{itemize}
\end{jebakan}

\begin{variasi}
Dari $\log2=0{,}301,\log3=0{,}477$: hitung $\log6=\log2+\log3=0{,}778$,
$\log 1{,}5=\log3-\log2=0{,}176$, $\log 8 = 3\log2=0{,}903$.
\end{variasi}

% ---------------- SOAL 1.3.8 ----------------
\soalno{1.3.8}{persamaan logaritma dengan argumen aljabar}
\begin{soalbox}{}
Selesaikan ${}^{3}\!\log(x-1) = 2$.
\end{soalbox}

\jawab
\langkah{1} Ubah ke bentuk eksponen (definisi): ${}^{3}\!\log(x-1)=2 \iff
x-1 = 3^2$.
\langkah{2} $3^2 = 9$, jadi $x-1=9 \Rightarrow x=10$.
\langkah{3} \textbf{Cek syarat} argumen log harus positif: $x-1>0 \Rightarrow
x>1$. Karena $x=10>1$, solusi \emph{valid}.
\jadi $x=10$.

\begin{jebakan}
\begin{itemize}[leftmargin=*,topsep=2pt]
  \item \textbf{Lupa cek domain ($x-1>0$).} Pada persamaan log \emph{selalu} cek
  argumen positif; solusi yang membuat argumen $\le0$ harus dibuang. (Di soal
  lain, ini bisa menggugurkan jawaban!)
  \item \textbf{Membalik: $x-1=2^3$.} Basisnya $3$, bukan $2$. Yang benar $3^2$.
\end{itemize}
\end{jebakan}

\begin{variasi}
${}^{2}\!\log(x+3)=4\Rightarrow x=13$; kasus yang menggugurkan:
${}^{2}\!\log(x^2-4)=0\Rightarrow x^2-4=1\Rightarrow x=\pm\sqrt5$ (cek keduanya
buat argumen positif). Atau dua log: ${}^{2}\!\log x+{}^{2}\!\log(x-2)=3$.
\end{variasi}

% ---------------- SOAL 1.3.9 ----------------
\soalno{1.3.9}{rantai tiga logaritma}
\begin{soalbox}{}
Nyatakan ${}^{a}\!\log b \cdot {}^{b}\!\log c \cdot {}^{c}\!\log a$ dalam bentuk
paling sederhana.
\end{soalbox}

\jawab
\langkah{1} Terapkan sifat rantai pada dua faktor pertama:
${}^{a}\!\log b\cdot{}^{b}\!\log c = {}^{a}\!\log c$.
\langkah{2} Kalikan dengan faktor ketiga (rantai lagi):
\[
{}^{a}\!\log c \cdot {}^{c}\!\log a = {}^{a}\!\log a = 1.
\]
\jadi hasilnya $\boxed{1}$.

\begin{tips}
\textbf{Cara pandang lain (ganti basis).} Tulis semua ke $\ln$:
$\dfrac{\ln b}{\ln a}\cdot\dfrac{\ln c}{\ln b}\cdot\dfrac{\ln a}{\ln c}$. Semua
faktor saling coret $\Rightarrow 1$. Visual ``coret berantai'' ini membantu
melihat mengapa hasilnya selalu $1$.
\end{tips}

\begin{jebakan}
\begin{itemize}[leftmargin=*,topsep=2pt]
  \item \textbf{Mengira hasilnya $abc$ atau $0$.} Rantai tertutup (kembali ke
  basis awal $a$) selalu $=1$.
\end{itemize}
\end{jebakan}

\begin{variasi}
${}^{2}\!\log3\cdot{}^{3}\!\log5\cdot{}^{5}\!\log2$ ($=1$), atau rantai tak
tertutup ${}^{a}\!\log b\cdot{}^{b}\!\log c={}^{a}\!\log c$ (tidak $=1$).
\end{variasi}

% ---------------- SOAL 1.3.10 ----------------
\soalno{1.3.10}{ganti basis dengan data log10}
\begin{soalbox}{}
Tentukan ${}^{2}\!\log 5$ jika $\log 2 = 0{,}301$ dan $\log 5 = 0{,}699$.
\end{soalbox}

\begin{ingat}{Sifat ganti basis (5)}
${}^{a}\!\log x = \dfrac{\log x}{\log a}$ (memakai log basis $10$ di ruas kanan).
Gunanya: mengubah log ``basis aneh'' menjadi log basis $10$ yang nilainya
diketahui/ada di tabel.
\end{ingat}

\jawab
\langkah{1} Ganti basis $2$ ke basis $10$:
\[
{}^{2}\!\log5 = \frac{\log5}{\log2}.
\]
\langkah{2} Masukkan nilai: $\dfrac{0{,}699}{0{,}301}$.
\langkah{3} Bagi: $\dfrac{0{,}699}{0{,}301}\approx 2{,}32$.
\jadi ${}^{2}\!\log5 \approx 2{,}32$. (Masuk akal: $2^2=4<5<8=2^3$, jadi
nilainya antara $2$ dan $3$.)

\begin{jebakan}
\begin{itemize}[leftmargin=*,topsep=2pt]
  \item \textbf{Terbalik jadi $\dfrac{\log2}{\log5}$.} Aturan: argumen di
  \emph{atas}, basis di \emph{bawah}. ${}^{2}\!\log5\Rightarrow\frac{\log5}{\log2}$.
  \item \textbf{Mengurangi, bukan membagi} ($0{,}699-0{,}301$). Ganti basis
  memakai \emph{pembagian}, bukan pengurangan.
  \item \textbf{Tidak menaksir kewajaran.} Selalu cek: $5$ ada di antara
  $2^2$ dan $2^3$, jadi jawab harus antara $2$ dan $3$.
\end{itemize}
\end{jebakan}

\begin{variasi}
${}^{5}\!\log2=\frac{\log2}{\log5}=\frac{0{,}301}{0{,}699}\approx0{,}43$
(kebalikannya!); ${}^{4}\!\log5=\frac{\log5}{\log4}=\frac{0{,}699}{2\cdot0{,}301}$.
Inilah tepat langkah yang menyelesaikan \emph{studi kasus}
($n={}^{1{,}5}\!\log100$).
\end{variasi}

% ============================================================
\begin{tips}
\textbf{Ringkasan strategi Bab 1.}
\begin{enumerate}[leftmargin=*,topsep=2pt]
  \item \emph{Eksponen:} ubah semua ke \textbf{basis prima} yang sama, lalu
  pakai 8 sifat. Pangkat negatif = kebalikan; pangkat pecahan = akar.
  \item \emph{Persamaan/pertidaksamaan:} samakan basis $\to$ samakan pangkat.
  Basis $<1$ $\Rightarrow$ balik tanda. Bentuk kuadrat $\to$ pemisalan $p=a^x>0$.
  \item \emph{Logaritma:} invers eksponen. Kali$\to$tambah, bagi$\to$kurang,
  pangkat$\to$kali; rantai untuk basis-tengah sama; ganti basis untuk memakai
  data $\log$ basis 10. \textbf{Selalu cek domain} pada persamaan log.
\end{enumerate}
\end{tips}

\end{document}
